\section{Conclusion and Future Work}
\label{sec:conclusion}

In this paper, we presented a rigorous empirical deconstruction of test-time adaptive model merging, revealing two critical vulnerabilities that have been overlooked in the literature. First, we exposed the \textbf{Overfitting-Optimizer Paradox} via a spatial shuffling diagnostic, demonstrating that standard, unconstrained layer-wise optimization does not capture stable layer-specific interactions, but instead overfits to local calibration statistics. Second, we identified the \textbf{Sacrificial Task Bias}, explaining how joint entropy minimization systematically sacrifices complex, high-entropy tasks (e.g., SVHN) to minimize joint loss.

To resolve these limitations, we introduced \textbf{RegCalMerge}, a robust calibration-aware test-time model merging framework comprising a core calibration engine (\textbf{CalMerge}) and an optional structural stabilizer (\textbf{ESR}). The framework incorporates:
\begin{itemize}
    \item \textbf{Elastic Spatial Regularization (ESR)}, which applies a dual Proximity Penalty ($\beta$) and Spatial Deviation Penalty ($\gamma$) to smooth coefficient variance and mitigate transductive overfitting;
    \item \textbf{Class-Capacity Normalization (CCN)}, which standardizes task-wise entropy values to a uniform $[0, 1]$ scale across different category structures; and
    \item \textbf{Scale-Normalized Entropy Weighting (SNEW)}, which balances the optimization landscape by weighting tasks based on their baseline entropy limits.
\end{itemize}

Through extensive evaluations over multiple random seeds and a dense 2D hyperparameter grid sweep ($\beta \times \gamma$), we demonstrated that our core calibration engine, CalMerge, successfully resolves the sacrificial task bias to achieve state-of-the-art Joint Mean accuracy (\textbf{61.82\%}), outperforming all standard baselines. Simultaneously, we showed that the optional ESR stabilizer allows researchers to navigate a smooth Generalization-Regularization Trade-off, providing a stable, highly predictable optimization surface that prevents parameter drift in ultra-noisy settings.

\subsection{Future Directions}
Our findings open several promising avenues for future model merging and test-time adaptation research:
\begin{enumerate}
    \item \textbf{Dynamic Test-Time Hyperparameter Tuning}: In our experiments, we showed that the generalization surface is exceptionally smooth under $\beta$ and $\gamma$. Future work could investigate metadata-driven or online controllers that dynamically adapt these penalties during the test-time stream based on incoming sample diversity or prediction confidence.
    \item \textbf{Cross-Architecture and Cross-Modal Adaptation}: While this work analyzed vision-language CLIP transformers, extending RegCalMerge to large language models (LLMs) or multimodal networks could help stabilize test-time adapter merging under severe distribution shifts.
    \item \textbf{Integration with Flatness-Aware Optimization}: Exploring the synergy between regularized test-time model merging and pre-training optimizers (such as Sharpness-Aware Minimization) may establish scaling laws that link weight-space flatness to test-time adaptation viability.
\end{enumerate}